% ========================== 第五节 业务与技术 ================================
% 主营业务与产品 → 经营模式 → 行业情况 → 客户/供应商 → 核心技术与研发
% → 知识产权 → 经营资质。收入构成、客户、供应商、研发数据与第二、六节勾稽。
% 样例业务：人工智能大模型（大模型解决方案 / 开放平台 MaaS / 智能应用）。
\gssection{业务与技术}

\subsection{发行人主营业务及主要产品}

\subsubsection{主营业务概述}

公司是一家专注于人工智能大模型的高新技术企业，主营业务为大模型的研发与产业
化应用。公司自主研发了通用大模型及多模态模型体系，围绕``模型研发—平台服务
—行业落地''构建业务闭环：面向行业客户提供大模型解决方案与私有化部署服务，
通过开放平台以API方式对外输出模型能力，并面向最终用户提供智能助手等智能应
用产品，帮助客户实现业务流程的智能化升级。

\subsubsection{主要产品及服务}

\paragraph{大模型解决方案}
面向金融、政务、制造、互联网等行业客户，基于公司自研基础模型提供行业模型联
合建模、私有化部署、智能体平台建设及配套系统集成服务，按项目制签约交付。

\paragraph{开放平台（MaaS）服务}
通过开放平台向企业客户及开发者提供基础模型与多模态模型的API调用、模型微调
与托管服务，按Token用量计量计费或以资源包方式预付费。

\paragraph{智能应用及其他}
面向最终用户的智能助手等应用产品的会员订阅与授权收入，以及技术咨询、培训等
其他服务。

\subsubsection{报告期主营业务收入构成}

报告期各期，公司主营业务收入按产品类别构成如下：

\biaodw{万元}
\begin{gstab}\zihao{-5}
\begin{longtable}{|L{27mm}|R{18mm}|R{14mm}|R{18mm}|R{14mm}|R{18mm}|R{14mm}|}
\hline
\thA{\multirow{2}{*}{产品类别}} & \multicolumn{2}{c|}{\bfseries 2025年度} & \multicolumn{2}{c|}{\bfseries 2024年度} & \multicolumn{2}{c|}{\bfseries 2023年度} \\ \cline{2-7}
\thA{} & \thB{金额} & \thB{占比} & \thB{金额} & \thB{占比} & \thB{金额} & \thB{占比} \\ \hline
\endfirsthead
\hline
\thA{\multirow{2}{*}{产品类别}} & \multicolumn{2}{c|}{\bfseries 2025年度} & \multicolumn{2}{c|}{\bfseries 2024年度} & \multicolumn{2}{c|}{\bfseries 2023年度} \\ \cline{2-7}
\thA{} & \thB{金额} & \thB{占比} & \thB{金额} & \thB{占比} & \thB{金额} & \thB{占比} \\ \hline
\endhead
大模型解决方案 & 72,845.36 & 55.21\% & 57,806.38 & 55.20\% & 44,731.24 & 54.52\% \\ \hline
开放平台（MaaS）服务 & 41,257.89 & 31.27\% & 33,193.42 & 31.70\% & 26,382.15 & 32.16\% \\ \hline
智能应用及其他 & 17,832.56 & 13.52\% & 13,714.05 & 13.10\% & 10,932.22 & 13.32\% \\ \hline
\multicolumn{1}{|c|}{\bfseries 主营业务收入} & {\bfseries 131,935.81} & {\bfseries 100.00\%} & {\bfseries 104,713.85} & {\bfseries 100.00\%} & {\bfseries 82,045.61} & {\bfseries 100.00\%} \\ \hline
\end{longtable}
\end{gstab}
\tabnote{报告期各期，公司另有少量技术咨询等其他业务收入，分别为410.71万元、
524.62万元和658.37万元。}

\subsection{发行人的经营模式}

\subsubsection{盈利模式}

公司通过大模型解决方案的项目制交付、开放平台服务的按量计费以及智能应用的订
阅授权获取收入与利润。开放平台服务按输入、输出Token用量计量计费或以资源包
预付费；智能应用按会员订阅周期收费；解决方案类业务按项目制组织实施，总体流
程如下：

% 业务流程图（式样对照 2026 年申报稿的工艺/业务流程图：方框+箭头链）
\begin{center}
\begin{tikzpicture}[x=1mm, y=1mm,
  flow/.style={draw=gsbar, semithick, fill=gsbarlight!35, minimum width=24mm,
    minimum height=9mm, align=center, font=\zihao{5}},
  arr/.style={-{Stealth[length=2.5mm]}, semithick, gsbar}]
\node[flow] (s1) at (0,0)   {合同签订};
\node[flow] (s2) at (30,0)  {方案设计};
\node[flow] (s3) at (60,0)  {模型适配\\与训练};
\node[flow] (s4) at (90,0)  {系统集成\\与部署};
\node[flow] (s5) at (120,0) {交付验收\\与运维};
\draw[arr] (s1) -- (s2); \draw[arr] (s2) -- (s3);
\draw[arr] (s3) -- (s4); \draw[arr] (s4) -- (s5);
\end{tikzpicture}
\end{center}
\gsfigcap{解决方案类业务实施流程}

\subsubsection{研发模式}

公司实行``基础模型—后训练对齐—产品化''三层研发体系：基础模型团队负责预训
练模型体系的迭代，对齐与安全团队负责监督微调、RLHF及安全防护能力建设，产品
团队面向行业场景进行智能体与应用开发；数据平台与训练平台为各层研发提供语料
与算力支撑。

\subsubsection{采购模式}

公司采购的主要内容包括算力资源（智算中心算力租赁及自建集群的服务器、智能计
算芯片等设备采购）、数据资源（语料采购与数据标注服务）以及云服务、基础软件
等，实行``合格供应商名录＋比价采购''的管理模式，对关键算力资源实施多源供应
策略。

\subsubsection{交付模式}

开放平台服务及智能应用通过云端在线交付；解决方案类业务根据客户数据安全要求，
以私有化部署或专属云方式交付，并按项目里程碑组织实施与验收。

\subsubsection{销售模式}

公司以直销模式为主，针对重点行业客户设立行业线团队，提供从方案设计到运维运
营的全流程服务；同时与独立软件开发商、系统集成商等生态伙伴合作拓展客户覆盖。

\subsection{发行人所处行业的基本情况}

\subsubsection{行业主管部门、监管体制及主要产业政策}

公司所处行业为人工智能行业，行业主管部门包括国家互联网信息办公室、工业和信
息化部、国家发展和改革委员会、科学技术部等。生成式人工智能服务适用《网络安
全法》《数据安全法》《个人信息保护法》《生成式人工智能服务管理暂行办法》
《互联网信息服务算法推荐管理规定》《互联网信息服务深度合成管理规定》等法律
法规及监管规则，实行算法备案与生成式人工智能服务备案等管理制度。《新一代人
工智能发展规划》及``人工智能＋''行动等政策明确支持人工智能核心技术攻关与产
业化应用，为行业发展提供了良好的政策环境。

\subsubsection{行业发展状况及市场规模}

大模型技术的突破显著拓展了人工智能的通用能力边界，推动人工智能从感知智能向
认知智能与生成智能演进。随着模型能力持续提升、推理成本下降及行业应用加速渗
透，大模型在办公、金融、政务、制造、互联网等领域的商业化落地不断深化。据公
开行业研究资料，我国人工智能核心产业规模已超过6,000亿元并保持快速增长，其
中大模型及生成式人工智能细分市场增速显著高于行业平均水平，市场规模情况如下：

% 市场规模柱状图（式样对照 2026 年申报稿：浅蓝=历史确定值、深蓝=预测值、
% 右上 CAGR 徽标、柱顶数值、E 为预测年份；数据为虚构示例）
\begin{center}
\begin{tikzpicture}[x=1mm, y=1mm]
\sffamily
% 标题与单位
\node[anchor=west, font=\heiti\zihao{5}] at (0,52) {中国人工智能大模型市场规模};
\node[anchor=west, font=\zihao{-5}] at (0,47) {单位：亿元};
% CAGR 徽标
\draw[gsbar, semithick] (95,46.5) rectangle (136,55.5);
\fill[gsbar] (95,46.5) rectangle (106,55.5);
\node[text=white, font=\zihao{-5}] at (100.5,51) {CAGR};
\node[anchor=west, font=\zihao{-5}] at (108,53.2) {2025---2028E};
\node[anchor=west, font=\heiti\zihao{-4}] at (108,49) {52.9\%};
% 柱体（历史浅蓝 2021-2025，预测深蓝 2026E-2028E；高度=数值/50）
\foreach \x/\v/\c in {0/12.5/gsbarlight, 17.5/38.6/gsbarlight, 35/147.2/gsbarlight,
  52.5/316.5/gsbarlight, 70/588.4/gsbarlight, 87.5/972.6/gsbar,
  105/1486.3/gsbar, 122.5/2105.8/gsbar}{
  \fill[\c] (\x,0) rectangle ({\x+12},{\v/50});
  \node[font=\zihao{-5}, anchor=south] at ({\x+6},{\v/50}) {\v};
}
% 年份轴
\foreach \x/\y in {0/2021, 17.5/2022, 35/2023, 52.5/2024, 70/2025,
  87.5/2026E, 105/2027E, 122.5/2028E}{
  \node[font=\zihao{-5}, anchor=north] at ({\x+6},-0.5) {\y};
}
\draw[semithick] (-1,0) -- (136,0);
% 图例
\fill[gsbarlight] (98,-7) rectangle (101,-9.2);
\node[anchor=west, font=\zihao{-5}] at (101.5,-8.1) {确定值};
\fill[gsbar] (115,-7) rectangle (118,-9.2);
\node[anchor=west, font=\zihao{-5}] at (118.5,-8.1) {预测值};
% 数据来源
\node[anchor=west, font=\zihao{-5}] at (0,-8.1) {数据来源：公开行业研究资料整理};
\end{tikzpicture}
\end{center}
\gsfigcap{中国人工智能大模型市场规模（示例数据）}

\subsubsection{行业竞争格局及发行人的行业地位}

行业竞争格局呈现``头部科技企业、专业大模型厂商与开源社区并存''的特点：头部
科技企业依托算力与生态优势布局通用大模型，专业厂商聚焦模型能力与行业落地，
开源社区则加速技术扩散。公司作为境内较早实现大模型商业化落地的专业厂商之一，
凭借自研模型体系、行业解决方案交付能力及开放平台生态，在金融、政务等重点行
业形成了领先的市场地位，与下游头部客户建立了稳定的合作关系。

\subsubsection{行业上下游关系}

行业上游主要为智能计算芯片、服务器与智算中心、云计算服务及数据资源供应行业，
其中高性能智能计算芯片的供应受国际贸易环境影响存在一定不确定性；下游覆盖金
融、政务、制造、互联网等众多行业，下游行业的数字化预算投入与智能化改造需求
直接影响本行业的市场空间。行业产业链结构如下：

% 产业链示意图（式样对照 2026 年申报稿：上/中/下游三层带状图）
\begin{center}
\begin{tikzpicture}[x=1mm, y=1mm,
  band/.style={draw=gsbar, semithick, minimum height=11mm, align=center,
    font=\zihao{5}},
  tag/.style={fill=gsbar, text=white, minimum width=16mm, minimum height=11mm,
    align=center, font=\heiti\zihao{5}},
  item/.style={draw=gsbar, semithick, fill=white, minimum width=29mm,
    minimum height=7mm, align=center, font=\zihao{-5}},
  arr/.style={-{Stealth[length=2.5mm]}, semithick, gsbar}]
% 上游
\node[tag] (t1) at (0,32) {上游};
\node[band, fill=gsbarlight!25, minimum width=126mm, anchor=west] (b1) at (9,32) {};
\node[item] at (26,32) {智能计算芯片};
\node[item] at (58,32) {服务器与智算中心};
\node[item] at (90,32) {云计算服务};
\node[item] at (122,32) {数据资源};
% 中游
\node[tag] (t2) at (0,16) {中游};
\node[band, fill=gsbarlight!45, minimum width=126mm, anchor=west] (b2) at (9,16) {};
\node[item, minimum width=45mm] at (34,16) {基础模型研发（发行人）};
\node[item, minimum width=45mm] at (106,16) {开放平台与模型服务};
% 下游
\node[tag] (t3) at (0,0) {下游};
\node[band, fill=gsbarlight!25, minimum width=126mm, anchor=west] (b3) at (9,0) {};
\node[item] at (26,0) {金融};
\node[item] at (58,0) {政务};
\node[item] at (90,0) {制造};
\node[item] at (122,0) {互联网及其他};
\draw[arr] (72,26.2) -- (72,21.8);
\draw[arr] (72,10.2) -- (72,5.8);
\end{tikzpicture}
\end{center}
\gsfigcap{人工智能大模型行业产业链}

\subsection{发行人的销售情况与主要客户}

报告期各期，公司向前五大客户的销售收入合计分别为35,090.50万元、45,547.75万
元和58,232.93万元，占营业收入的比例分别为42.56\%、43.28\%和43.92\%。公司与
主要客户均无关联关系，不存在向单个客户的销售比例超过总额50\%或严重依赖于少
数客户的情形。2025年度公司前五大客户销售情况如下：

\biaodw{万元}
\begin{gstab}
\begin{longtable}{|C{12mm}|L{40mm}|R{34mm}|R{34mm}|}
\hline
\thA{序号} & \thB{客户名称} & \thB{销售金额} & \thB{占营业收入比例} \\ \hline
\endfirsthead
\hline
\thA{序号} & \thB{客户名称} & \thB{销售金额} & \thB{占营业收入比例} \\ \hline
\endhead
1 & 客户A & 18,563.19 & 14.00\% \\ \hline
2 & 客户B & 12,845.67 & 9.69\% \\ \hline
3 & 客户C & 10,236.45 & 7.72\% \\ \hline
4 & 客户D & 8,952.34 & 6.75\% \\ \hline
5 & 客户E & 7,635.28 & 5.76\% \\ \hline
\multicolumn{2}{|c|}{\bfseries 合计} & {\bfseries 58,232.93} & {\bfseries 43.92\%} \\ \hline
\end{longtable}
\end{gstab}
\tabnote{同一控制下的客户已合并计算；2023年度、2024年度前五大客户明细在正式
申报文件中应按相同格式逐年披露。}

\subsection{发行人的采购情况与主要供应商}

报告期各期，公司算力资源、数据资源及软硬件等采购总额分别为44,586.24万元、
55,238.65万元和68,435.26万元，向前五大供应商的采购金额合计分别为20,218.62万元、
25,121.85万元和31,176.59万元，占当期采购总额的比例分别为45.35\%、45.48\%和
45.56\%，主要供应商为智算中心、云计算服务商及服务器设备厂商。公司与主要供
应商均无关联关系，不存在向单个供应商的采购比例超过总额50\%或严重依赖于少数
供应商的情形。2025年度公司前五大供应商采购情况如下：

\biaodw{万元}
\begin{gstab}
\begin{longtable}{|C{12mm}|L{40mm}|R{34mm}|R{34mm}|}
\hline
\thA{序号} & \thB{供应商名称} & \thB{采购金额} & \thB{占采购总额比例} \\ \hline
\endfirsthead
\hline
\thA{序号} & \thB{供应商名称} & \thB{采购金额} & \thB{占采购总额比例} \\ \hline
\endhead
1 & 供应商A & 8,935.24 & 13.06\% \\ \hline
2 & 供应商B & 7,246.85 & 10.59\% \\ \hline
3 & 供应商C & 5,832.47 & 8.52\% \\ \hline
4 & 供应商D & 4,926.35 & 7.20\% \\ \hline
5 & 供应商E & 4,235.68 & 6.19\% \\ \hline
\multicolumn{2}{|c|}{\bfseries 合计} & {\bfseries 31,176.59} & {\bfseries 45.56\%} \\ \hline
\end{longtable}
\end{gstab}

\subsection{核心技术与研发情况}

\subsubsection{核心技术情况}

公司围绕大模型训练与推理两大技术主线，形成了六项核心技术，均来源于自主研发，
且均已应用于主要产品及服务并形成收入：

\begin{gstab}
\begin{longtable}{|C{8mm}|L{42mm}|L{56mm}|C{18mm}|}
\hline
\thA{序号} & \thB{核心技术名称} & \thB{主要应用} & \thB{技术来源} \\ \hline
\endfirsthead
\hline
\thA{序号} & \thB{核心技术名称} & \thB{主要应用} & \thB{技术来源} \\ \hline
\endhead
1 & 大规模分布式训练与算力调度技术 & 千卡级集群高效并行训练、故障容错与算力利用率优化 & 自主研发 \\ \hline
2 & 混合专家（MoE）模型架构技术 & 在可控推理成本下扩展模型参数规模与模型能力 & 自主研发 \\ \hline
3 & 多模态理解与生成技术 & 文本、图像、音频、视频的统一表征、理解与生成 & 自主研发 \\ \hline
4 & 模型对齐与安全防护技术 & 基于RLHF的偏好对齐、内容安全过滤与幻觉抑制 & 自主研发 \\ \hline
5 & 推理加速与模型压缩技术 & 量化、蒸馏与稀疏化，降低推理时延与单位算力成本 & 自主研发 \\ \hline
6 & 智能体编排与工具调用技术 & 任务规划、多智能体协同及企业系统工具调用 & 自主研发 \\ \hline
\end{longtable}
\end{gstab}

\subsubsection{研发组织架构}

公司按``基础研究—平台工程—产品化''设置研发组织，主要研发部门及职责如下：

\begin{gstab}
\begin{longtable}{|C{10mm}|L{30mm}|L{92mm}|}
\hline
\thA{序号} & \thB{部门} & \thB{主要职责} \\ \hline
\endfirsthead
\hline
\thA{序号} & \thB{部门} & \thB{主要职责} \\ \hline
\endhead
1 & 基础模型研发部 & 负责基础模型架构研究、大规模预训练与模型能力迭代 \\ \hline
2 & 多模态研发部 & 负责图像、音频、视频等多模态理解与生成模型的研发 \\ \hline
3 & 对齐与安全部 & 负责监督微调、基于人类反馈的强化学习及内容安全防护体系建设 \\ \hline
4 & 智能体与应用平台部 & 负责智能体编排平台、开放平台及行业应用产品的研发 \\ \hline
5 & 数据平台部 & 负责训练语料建设、数据清洗标注体系与数据合规管理 \\ \hline
6 & 算力工程部 & 负责智算集群建设运维、分布式训练框架与推理加速优化 \\ \hline
\end{longtable}
\end{gstab}

\subsubsection{研发投入情况}

报告期各期，公司研发投入及占营业收入的比例如下：

\biaodw{万元}
\begin{gstab}
\begin{longtable}{|L{49mm}|R{29mm}|R{29mm}|R{29mm}|}
\hline
\thA{项目} & \thB{2025年度} & \thB{2024年度} & \thB{2023年度} \\ \hline
\endfirsthead
\hline
\thA{项目} & \thB{2025年度} & \thB{2024年度} & \thB{2023年度} \\ \hline
\endhead
研发投入金额 & 12,890.73 & 9,825.64 & 7,412.35 \\ \hline
占营业收入的比例 & 9.72\% & 9.34\% & 8.99\% \\ \hline
\end{longtable}
\end{gstab}
\tabnote{研发投入为费用化的研发费用；用于模型训练的算力集群等资本性投入未计
入上表。}

报告期内，公司累计研发投入30,128.72万元，占累计营业收入的比例为9.41\%。截
至2025年12月31日，公司研发与技术人员726人，占员工总数的56.45\%，符合科创板
对科创属性的相关要求。

\subsubsection{知识产权情况}

截至2025年12月31日，公司及子公司拥有的知识产权情况如下：

\begin{gstab}
\begin{longtable}{|L{40mm}|R{30mm}|L{56mm}|}
\hline
\thA{类别} & \thB{数量（项）} & \thB{备注} \\ \hline
\endfirsthead
\hline
\thA{类别} & \thB{数量（项）} & \thB{备注} \\ \hline
\endhead
发明专利 & 186 & 均与主营业务相关，其中应用于核心技术的发明专利118项 \\ \hline
实用新型专利 & 32 & —— \\ \hline
外观设计专利 & 18 & —— \\ \hline
软件著作权 & 120 & 涵盖模型训练平台、推理引擎、开放平台及智能体平台等 \\ \hline
\end{longtable}
\end{gstab}

公司核心技术相关专利均系原始取得，不存在与他人共有或存在权属纠纷的情形。

\subsubsection{在研项目情况}

截至2025年12月31日，公司主要在研项目包括新一代多模态基础模型、长上下文与记
忆机制优化、端侧轻量化模型、行业智能体平台及大模型安全评测体系等，均围绕现
有主业展开。

\subsubsection{核心技术人员}

公司核心技术人员为卫××、蒋××、沈××、韩××4名，均于公司设立当年或次年
加入公司，主持或参与了公司全部核心技术的研发。报告期内，公司核心技术人员未
发生变动。

\subsection{主要经营资质}

公司为高新技术企业，已按规定完成生成式人工智能服务备案及深度合成服务算法备
案，取得增值电信业务经营许可证，并通过ISO/IEC~27001信息安全管理体系认证和
信息系统安全等级保护三级测评。公司已取得开展现有业务所需的全部资质与认证，
报告期内不存在因资质瑕疵受到行政处罚的情形。
